\documentclass{article}
\usepackage[left=2cm, right=2cm, top=2cm, bottom=2cm]{geometry}
\usepackage{graphicx}
\usepackage{amsmath}
\usepackage{amssymb}
\usepackage{amsthm}
\usepackage{fancyhdr}
\usepackage{verbatim}
\usepackage{listings}
\usepackage{xcolor}
\usepackage{pdfpages}
\usepackage{cancel}
\usepackage{physics}
\usepackage{esint}
\usepackage{braket}

\lstset{
    language=C++,
    basicstyle=\ttfamily\footnotesize,
    keywordstyle=\color{blue},
    commentstyle=\color{green},
    stringstyle=\color{red},
    numbers=left,
    numberstyle=\tiny\color{gray},
    frame=single,
    breaklines=true,
    captionpos=b,
}


\title{MTH 446: Lecture \# 13}
\author{Cliff Sun}

\newtheorem{theorem}{Theorem}[section]
\newtheorem{lemma}[theorem]{Lemma}
\newtheorem{definition}[theorem]{Definition}
\newtheorem{conjecture}[theorem]{Conjecture}
\newtheorem{proposition}[theorem]{Proposition}
\newtheorem{corollary}[theorem]{Corollary}
\newtheorem{one minute paper}[theorem]{One Minute Paper}

\pagestyle{fancy}
\lhead{\textbf{\thepage}\ \ \nouppercase{\rightmark}}
\chead{MTH 446: Lecture \# 13}
\rhead{Cliff Sun}

\begin{document}

\maketitle

\section*{Lecture Span}
\begin{enumerate}
    \item Enter here
\end{enumerate}

\section*{Uniqueness property}

\begin{definition}
    If $\{z \in D, f(z) = g(z)\}$, then $f=g$ for all $z \in D$. Note, $f$ and $g$ must be analytic functions on $\mathbb{C}$. 
\end{definition}

\subsection*{Example}

Note that $\sin z$ and $\cos z$ are analytic in $\mathbb{C}$. Then, we wish to prove that 

\begin{equation}
    \cos^2 z+ \sin^2 z= 1
\end{equation}

Since we know that $\cos^2 x + \sin^2 x = 1$ for all $x\in \mathbb{R}$. Therefore, since $\cos,\sin \in H(\mathbb{C})$, then it follows that 

\begin{equation}
    f(z) = \cos^2z + \sin^2 z = g(z) = 1, \quad \forall z \in \mathbb{C}
\end{equation}

Taylor's expansion theorem, if $f \in H(D_R(z_0))$, then 

\begin{equation}
    f(z) = \sum a_n (z - z_0)^n, \quad |z-z_0| < R
\end{equation}

Note that such a function is infinitely differentiable. Analytic functions are relative to polynomials.  

\begin{theorem}
    Suppose $f'(z) = 0$ for all $z \in D$. Then, $f$ is constant everywhere in $D$.  
\end{theorem}

\begin{proof}
    Let $f = u + iv$. Then, by the CR equations, we have 
    \begin{equation}
        u_x = v_y, \quad u_y = -v_x
    \end{equation}
    Therefore, 
    \begin{equation}
        f'(z) = u_x + iv_x = 0
    \end{equation}
    Then 
    \begin{align}
        u_x &\equiv 0 \equiv v_y \\
        v_x &\equiv 0 \equiv u_y
    \end{align}
    Therefore, 
    \begin{align}
        \nabla u \equiv 0 &\implies u = \text{const.}\\
        \nabla v \equiv 0 &\implies v = \text{const.}
    \end{align}
    Therefore, $f$ is constant throughout $D$. 
\end{proof}

Some notes:
\begin{enumerate}
    \item A point in which $f$ is not analytic is a singular point. 
    \item A singular point $z_0$ of $f$ is said to be isolated if $\exists \epsilon$, $f \in H(\overset{\circ}{D}_{\epsilon}(z_0))$.  
\end{enumerate}

\subsection*{Example}

\begin{equation}
    f(z) = \frac{1}{z^4}, \quad z \neq 0
\end{equation}

Note that $1 \in H(\mathbb{C})$ and $z^4 \in H(\mathbb{C})$. And $z^4 \neq 0$ if $z = 0$. Therefore, $f$ is origin across $\mathbb{C}$ expect for the origin. In other words, 

\begin{equation}
    f \in H(\mathbb{C}/\{0\})
\end{equation}

Therefore, $0$ is an isolated singular point. 

\subsection*{Example}

\begin{equation}
    f = \frac{2z + 1}{z(z^2 + 1)}
\end{equation}

The denominator is zero if $z = 0, i, -i$. Therefore, 

\begin{equation}
    f \in H(\mathbb{C}/\{0, i, -i\})
\end{equation}

Therefore, $0, i, -i$ are isolated singular points. Note, if the set is finite, then all the values in the set are isolated singular points. 

\subsection*{Example}

\begin{equation}
    f(z) = x^2 + y^2
\end{equation}

Both $u$ and $v$ are continuously differentiable on $\mathbb{C}$. 

Previously, if $z \neq 0$, then $f$ is not differentiable. However, since $f'$ is not differentiable on a neighborhood, it follows that $f'$ is not analytic at all. 

\subsection*{Example}

\begin{equation}
    f = xy + iy
\end{equation}

We have established that $f'$ DNE if $z \neq 2$ and $f'(2) = 0$. Therefore, we can establish that $f$ is not analytic at $z \neq 2$. Now, consider $z = 2$, no neighborhood. Therefore, $f$ is nowhere analytic. 

\begin{proposition}
    $f$ must be constant in $D$ if one of the following is satisfied. 
    \begin{enumerate}
        \item $f$ is a real-valued function 
        \item $\overline{f} \in H(D)$
        \item $|f(z)| = \text{const.}$ in $D$
    \end{enumerate}
\end{proposition}

\begin{proof}
    We first prove (a). Note that $f \in H(D)$. Also consider $f = u + i(0)$. Then, 
    \begin{equation}
        u_x = v_y = 0 \iff u = \text{const.}
    \end{equation}
    We can show that $\nabla u = 0$ and $\nabla v = 0$. Therefore, $f$ is constant on $D$.
    
    We next prove (b). We note that $f, \overline{f} \in H$. Note that we can rewrite $v$ and $u$ in terms of $f$ and $\overline{f}$. But then, $u$ and $v$ 
    are real valued. Therefore, by (a), we have that $f$ is const. in $D$.

    Finally, we prove (c). Note that $|f|^2 = c^2$. Also, $f \in D$. Firstly, if $c = 0$, then $f$ is 0 and therefore is constant. 
    Next, if $c \neq 0$, then note that $f$ can never be $0$ for all $z \in D$. Next, $\overline{f}f = c^2$. Then, 
    \begin{equation}
        \overline{f} = \frac{c^2}{f} \in H(D)
    \end{equation}
    Therefore, $f$ is constant in $D$ by (b). 
\end{proof}

We introduce notation: 
\begin{enumerate}
    \item $c^1(D)$ is the set of all real valued functions $u$ such that $u$, $u_x$, $u_y$ exist and are cont. in $D$
    \item $c^2(D)$ is the set of all real valued funcs $u$ such that $u$, $u_x$, $u_y$, $u_{xx}$, etc. exist and are cont. in $D$. 
\end{enumerate}

\begin{definition}
    \begin{enumerate}
        \item $c^0(D)$ is the space of continous functions in $D$
        \item $c^k(D)$ is the space of all functions $u$ with continuous partial derivatives of order $j$. 
        \item $c^{\infty}(D)$ is the set of all infinitely differentiable functions in $D$. 
    \end{enumerate}
\end{definition}

\begin{definition}
    A function $u$ is called harmonic if $u \in C^2$ and $\nabla^2 u = 0$. 
\end{definition}

Notation: Har(D) is the set of all harmonic functions in $D$. 

\begin{theorem}
    $f \in H(D)$ and $f = u + iv$. Then $u,v \in \text{Har}(D)$. 
\end{theorem}




\end{document}